\chapter{Introduction}
\label{ch:intro}

nornRTOS is a real time operating system for small microcontrollers. It consists of a
kernel, \emph{nornCore}, a TCP/IP stack, \emph{nornNET}, and a driver for the ESP32's
WiFi radio that runs without ESP-IDF or FreeRTOS, \emph{nornWIFI}. The kernel runs on the
8\,bit AVR family (ATmega328P on the Arduino UNO, ATmega2560 on the Arduino Mega) and on
the Xtensa LX6 of the ESP32; a host port runs the same kernel on Linux for testing. The
whole system is written in C99 with a few dozen lines of assembler per port, and every
piece of it is compiled from source together with the application and its configuration.

This guide is for the engineer who writes an application on nornRTOS. It explains the
model of the kernel (chapter~\ref{ch:concepts}), how an application is built and
configured (chapters~\ref{ch:start} and~\ref{ch:config}), how the primitives are meant to
be used (chapter~\ref{ch:primitives}) and what the kernel does when something goes wrong
(chapter~\ref{ch:robust}). Chapter~\ref{ch:api} is the reference of the public API,
chapter~\ref{ch:net} covers the network modules, and chapters~\ref{ch:perf}
and~\ref{ch:verify} present the measurements and the test suite the figures in this guide
rest on. The form of the guide follows the RTuinOS user guide~\cite{rtuinos}, the
operating system nornRTOS grew out of and replaced.

\section{What kind of kernel this is}

nornCore is a \emph{static}, \emph{fixed priority}, \emph{preemptive} kernel with
\emph{one task per priority}:

\begin{itemize}
\item \textbf{Static.} Every task is created in \code{main()} before the scheduler starts
  and lives for ever. The kernel never allocates memory: the application owns every stack
  and every synchronisation object. Nothing in the kernel grows at run time, so the
  memory an application needs is known when it is linked.
\item \textbf{Fixed priority, one task per priority.} A task's identifier \emph{is} its
  priority, 1 to \code{RTOS\_CFG\_NUM\_TASKS}; 0 is the idle task. Two tasks never share
  a priority, so there is no round robin, no time slice and no question about the order
  in which equal tasks run. The theory behind fixed priority scheduling and its
  schedulability tests is Liu and Layland's~\cite{liulayland} and its descendants, and
  nornCore gives the application exactly the model that theory assumes.
\item \textbf{Preemptive, $O(1)$.} The scheduler is one bit mask of ready tasks and one
  highest-bit lookup: between four and twenty cycles depending on the CPU. A task that
  becomes ready with a higher priority than the running one runs at once; a task that
  blocks hands the CPU to the highest ready one. Every blocking call has a timeout.
\item \textbf{Priority ceiling mutexes.} Mutual exclusion uses the immediate priority
  ceiling protocol~\cite{shaetal,baker}: a task that locks a mutex runs at the mutex's
  ceiling while it holds it, so with a correct ceiling nobody ever blocks on a mutex,
  priority inversion is bounded by one critical section, and deadlock is impossible.
\item \textbf{Measured.} The cost of every kernel call, the context switch, the interrupt
  latency and the stack each task needs are measured in CPU cycles on the real boards and
  in a cycle exact simulator, before and after every change to the kernel
  (chapter~\ref{ch:perf}). A kernel of a few kilobytes has no excuse for unknown numbers.
\end{itemize}

What is deliberately absent is as telling: time slicing, dynamic creation and deletion of
tasks at run time (the life cycle in section~\ref{sec:lifecycle} and
appendix~\ref{app:roadmap} are the preparation), priority inheritance (the ceiling
protocol is simpler and has better bounds) and symmetric multiprocessing (the two ESP32
cores run two independent kernels joined by a message link instead).

\section{The modules and the boards}

\begin{table}[htbp]
\centering\small
\begin{tabular}{@{}lll@{}}
\toprule
Module & What it is & Boards \\
\midrule
nornCore & the kernel, the ports, a small HAL (UART, GPIO, delay, trace) & UNO, Mega 2560, ESP32, host \\
nornNET  & Ethernet, ARP, IPv4, ICMP, UDP, TCP, DHCP client, DNS resolver & ESP32, Mega 2560 (SLIP), host \\
nornWIFI & the ESP32 radio as a station, WPA2 personal, a nornNET device & ESP32 \\
nornHTTP & an HTTP/1.1 client over any byte stream, no heap, streamed bodies & any, with a nornNET adapter \\
\bottomrule
\end{tabular}
\caption{The modules of nornRTOS.}
\end{table}

\begin{table}[htbp]
\centering\small
\begin{tabular}{@{}llllll@{}}
\toprule
Board & CPU & Clock & RAM & Flash & Simulator \\
\midrule
Arduino UNO      & ATmega328P, 8 bit AVR     & 16 MHz  & 2 KB   & 32 KB   & simavr (cycle exact) \\
Arduino Mega 2560 & ATmega2560, 8 bit AVR    & 16 MHz  & 8 KB   & 256 KB  & simavr (cycle exact) \\
ESP32 devkit     & ESP32-D0WDQ6, 2 Xtensa LX6 & 240 MHz & 520 KB & 4 MB    & Espressif QEMU (not cycle exact) \\
\bottomrule
\end{tabular}
\caption{The boards. The ESP32 runs bare metal: the ROM loads the image, no ESP-IDF.}
\label{tab:boards}
\end{table}

The UNO has room for the kernel and a small application; nornNET runs on the Mega over
a serial line and on the ESP32 over Ethernet (QEMU) or WiFi (nornWIFI). Everything in this
guide that gives a number says which board it was measured on.

\section{Conventions}

\begin{description}
\item[Checked] means a mandatory check, \code{RTOS\_CHECK}, compiled into every
  configuration. A violation stops the system with \code{RTOS\_PANIC\_ARGUMENT} and a
  record naming the values. Task ids, NULL storage, stacks below the minimum, timeouts
  above the limit and foreign blocks freed are checked.
\item[Asserted] means \code{RTOS\_ASSERT}, compiled only with \code{RTOS\_CFG\_ASSERT}.
  Without it a broken asserted contract is undefined behaviour. Recursive locks, unlock
  order and empty event masks are asserted.
\item[Task context] is a task function or anything it calls. \textbf{ISR body} is the
  code inside \code{RTOS\_ISR}, a software timer callback running in the tick, or the
  trace hook. The \textbf{idle task} is the context of \code{main()} after
  \code{rtos\_start()} and of the idle hook. Which functions each context may call is
  the subject of section~\ref{sec:contexts}.
\item[Ticks] are the kernel's unit of time, \code{RTOS\_CFG\_TICK\_HZ} per second
  (1\,kHz in every configuration shipped). \code{RTOS\_MS(ms)} converts milliseconds.
\item[Cycles] are CPU clock cycles: 62.5\,ns on the 16\,MHz AVR, 4.17\,ns on the
  240\,MHz ESP32. Most figures are in cycles because that is what the benchmark measures
  and what does not change with the clock.
\end{description}

Function names are written as in the header: \code{rtos\_semGive}. A function with the
suffix \code{FromIsr} is the variant for an ISR body. Names from the configuration start
with \code{RTOS\_CFG\_}.
